(1) 在两个界面反射振幅分别 \(r_{01}=(n_0-n_f)/(n_0+n_f)\)、\(r_{12}=(n_f-n_s)/(n_f+n_s)\)。将薄膜内的多次反射场以往返相位 \(e^{2i\beta}\) 求几何级数，可得总反射 \(r=(r_{01}+r_{12}e^{2i\beta})/(1+r_{01}r_{12}e^{2i\beta})\)，反射率 \(R=|r|^2\)。不可仅相加两次反射强度，因为不同光程的电场需要相干叠加。

(2) 设计四分之一波长条件为 \(n_fd=\lambda_0/4\)，即 \(\beta=\pi/2\)，往返因子为 \(-1\)，于是 \(r=(r_{01}-r_{12})/(1-r_{01}r_{12})\)。零反射要求 \(r_{01}=r_{12}\)，解两界面 Fresnel 比例得 \(n_f^2=n_0n_s\)，所以最优无耗散膜折射率是 \(n_f=\sqrt{n_0n_s}\)。这是单层正入射严格消反射条件，依赖设计波长与材料色散。

(3) 若 \(n_f=n_0\)，首界面反射为零，总反射幅度退化为 \(r_{12}e^{2i\beta}\)，反射率为裸介质—基底值 \([(n_0-n_s)/(n_0+n_s)]^2\)，相位变化不影响强度。若 \(n_f=n_s\)，末界面无反射，仍得到相同裸界面反射率。三层体系无吸收且基底无限延展，合适的能流归一化下反射和透射功率之和为一，零反射时全部入射能量进入基底。

(4) 光波长偏离 \(\lambda_0\) 时 \(\beta=(\pi/2)(\lambda_0/\lambda)\) 不再是四分之一波长，两个主要反射场不能完全相消。若膜厚偏差 \(\delta d\)，往返相位偏差约为 \(\delta(2\beta)=4\pi n_f\delta d/\lambda_0\)，导致非零残余反射；折射率偏差破坏 \(n_f^2=n_0n_s\)，即使相位正确仍不能完全消除反射。宽带增透需多层或梯度折射率设计，不能外推单层的单频匹配结论。
